% Chapter 3: machine overview.
\chapter{Machine overview}\label{ch:overview}
\index{machine overview}

The paper winder pulls paper off a supply, past a rubber measuring wheel,
and onto a roll. A small controller board counts how far the measuring wheel
turns, switches the motor on and off through a relay, and shows what it is
doing on a four-line screen. You start and stop winds with four buttons on
the panel.

\begin{figure}[h]
\centering
\begin{tikzpicture}[font=\small,>=Stealth,
  lab/.style={draw,fill=white,rounded corners=1pt,inner sep=3pt,align=center}]
  % paper supply
  \draw[line width=1pt] (0,0) circle (1.0);
  \draw (0,0) circle (0.25);
  \node[lab] (sup) at (0,-1.8) {Paper supply};
  % paper path
  \draw[line width=1.5pt] (0,1.0) -- (4.2,1.0);
  \draw[line width=1.5pt] (4.2,1.0) -- (8.0,1.0);
  % measuring wheel above the paper
  \draw[line width=1pt] (4.2,1.55) circle (0.55);
  \fill (4.2,1.55) circle (0.05);
  \draw[->] (4.2,1.55) ++(30:0.75) arc (30:150:0.75);
  \node[lab] (whl) at (4.2,3.2) {Measuring wheel\\with encoder\\(about 4 in across)};
  \draw (whl.south) -- (4.2,2.1);
  % take-up roll
  \draw[line width=1pt] (9.0,0.0) circle (1.0);
  \draw (9.0,0.0) circle (0.25);
  \draw[line width=1.5pt] (8.0,1.0) -- (9.0,1.0);
  \node[lab] (roll) at (9.0,-1.8) {Take-up roll};
  % motor
  \draw[line width=1pt] (10.6,-0.5) rectangle (11.8,0.5);
  \node at (11.2,0) {M};
  \draw[line width=1pt] (9.25,0) -- (10.6,0);
  \node[lab] (mot) at (11.2,1.6) {Motor};
  \draw (mot.south) -- (11.2,0.5);
  % control box
  \draw[line width=1.2pt] (3.0,-4.6) rectangle (8.6,-2.6);
  \node[anchor=north west,font=\small\bfseries] at (3.05,-2.65) {Control box};
  \draw (3.4,-3.2) rectangle (6.0,-4.2);
  \node[font=\scriptsize\ttfamily,align=left] at (4.7,-3.7) {LCD 20 x 4};
  \foreach \i/\x in {1/6.4,2/6.9,3/7.4,4/7.9} {
    \draw (\x,-4.0) circle (0.18);
    \node[font=\scriptsize] at (\x,-4.0) {\i};
    \fill (\x,-3.35) circle (0.07);
  }
  \node[font=\scriptsize] at (7.15,-3.05) {LEDs 1-4};
  \node[font=\scriptsize] at (7.15,-4.4) {Buttons 1-4};
  % relay box
  \draw[line width=1pt] (9.6,-3.9) rectangle (11.4,-3.0);
  \node[align=center,font=\small] at (10.5,-3.45) {Relay};
  \draw[->,dashed] (8.6,-3.45) -- (9.6,-3.45);
  \draw[->,dashed] (10.5,-3.0) -- (11.2,-0.5);
  \draw[->,dashed] (4.2,0.95) .. controls (3.0,-1.0) .. (3.6,-2.6);
  \node[font=\scriptsize,align=center] at (2.2,-1.3) {encoder\\signal};
  % serial
  \draw[dashed] (3.0,-4.3) -- (1.2,-4.3);
  \node[lab,font=\scriptsize] at (0.6,-4.3) {USB serial\\adapter\\to PC};
  \node[font=\footnotesize\itshape] at (5.8,-5.1) {Schematic, not to scale.};
\end{tikzpicture}
\caption{The main parts of the paper winder.}
\label{fig:overview}
\end{figure}

\section*{The main parts}
\begin{description}[style=nextline,leftmargin=1.5em]
  \item[Measuring wheel and encoder]\index{measuring wheel}\index{encoder}
    A rubber-faced wheel about 4 inches across rides on the paper. A
    rotation sensor (the encoder) on its shaft gives 2400 counts for every
    turn of the wheel, which is about 191 counts per inch of paper. This is
    the only thing the machine uses to measure length, so the wheel must be
    clean and must turn freely with the paper without slipping.
  \item[Motor and relay]\index{relay}\index{motor}
    The controller switches the motor on and off through a relay. It is
    simply on or off; there is no speed control. When the relay opens, the
    roll keeps turning for a short distance (the \emph{coast}).
  \item[Control box]
    Holds the controller board, the 20-character by 4-line screen, the four
    LEDs, the four buttons and a buzzer.
  \item[Serial connection]\index{serial adapter}
    Three wires (TX, RX, GND) from a USB serial adapter let a PC run the
    \textbf{winder} program for calibration and settings. The PC is not needed
    to wind paper.
  \item[USB port on the controller]
    Used only to load new firmware (Chapter~\ref{ch:flashing}). It does not
    talk to the winder program.
\end{description}

\begin{note}
Your machine's exact layout (which side the supply is on, where the power
switch is) may differ from the drawing. The power switch location is:
[POWER SWITCH LOCATION].
\end{note}
